\documentclass{article}
\usepackage[margin=1in]{geometry}
\usepackage{hyperref}
\usepackage{ulem}
\hypersetup{
    colorlinks=true,
    linkcolor=black,
    filecolor=magenta,
    urlcolor=cyan,
}
\urlstyle{same}

\title{Mistake Reference -- Writing Practice}
\author{Junzhang Luo}
\date{}

\newcommand{\mistakeentry}[6]{
\item \textbf{Frequent error:} #1\\
\textbf{Correct form:} #2\\
\textbf{Short rule:} #3\\
\textbf{Corrected example:} #4\\
\textbf{Frequency:} #5\\
\textbf{Sources:} #6
}

\begin{document}

\maketitle
\tableofcontents

\section{Methodology}
\begin{itemize}
    \item Primary source: \texttt{DELF/writing/practice.tex}.
    \item Errors are extracted from your \textbf{original drafts only}.
    \item \textit{Fixed/Correction/Corrected/Fixes} sections are used to infer target forms.
    \item \textit{B2/C1} versions are \textbf{not} treated as new error sources.
\end{itemize}

\section{Top Priorities (Spaced Review)}
\subsection*{Highest-frequency categories}
\begin{enumerate}
    \item Prepositions and articles (about 28 occurrences)
    \item Agreement: gender, number, adjectives, participles (about 24 occurrences)
    \item Verb forms and tense control (about 20 occurrences)
    \item Spelling and accents (about 18 occurrences)
\end{enumerate}

\subsection*{8 patterns to prioritize}
\begin{enumerate}
    \item \textit{de/du/des} after negation and quantity.
    \item \textit{à/en/au/aux} for places and countries.
    \item \textit{qui/que/dont} by grammatical function.
    \item Adjective-noun and subject-verb agreement.
    \item Correct auxiliary choice in \textit{passé composé} (\textit{être} vs \textit{avoir}).
    \item Word order with adverbs (\textit{a beaucoup changé}).
    \item High-value collocations (\textit{présenter un spectacle}, \textit{être surpris de voir que}).
    \item Connector accuracy (\textit{car/parce que}), avoid double concession markers.
\end{enumerate}

\subsection*{Suggested review rhythm}
\begin{itemize}
    \item Day 0: Read all 8 patterns + do every micro-drill.
    \item Day 2: Review categories 1--4 only.
    \item Day 5: Review all categories, focusing on persistent errors.
    \item Day 10: Write a short paragraph and self-correct with this sheet.
\end{itemize}

\section{Spelling and Accents}
\subsection*{Entries}
\begin{enumerate}
    \mistakeentry
    {\textit{j'éspiere}, \textit{recénment}, \textit{trés}, \textit{apartement}}
    {\textit{j'espère}, \textit{récemment}, \textit{très}, \textit{appartement}}
    {Accent marks and standard spelling are mandatory in formal writing.}
    {\textit{J'espère que tu vas bien et je cherche un nouvel appartement.}}
    {14 occurrences}
    {A1 Sample 2/3/4; A2 Sample 1 Ex.1; A2 Sample 2 Ex.1; A2 Sample 3 Ex.2; B1 Exemple 1; Writing Practice Day 1-3}

    \mistakeentry
    {Missing/swapped consonants: \textit{caompagne}, \textit{vinngt derinères}, \textit{environmentaux}}
    {\textit{campagne}, \textit{vingt dernières}, \textit{environnementaux}}
    {Double-check common consonant clusters (\textit{-nn-, -mm-, -rt-, -nt-}).}
    {\textit{Au cours des vingt dernières années, les changements environnementaux...}}
    {11 occurrences}
    {A1 Sample 3; B1 Exemple 1/2; A2 Sample 4 Ex.1; Writing Practice Day 1-3}

    \mistakeentry
    {Inconsistent capitalization: \textit{Le Samedi}, \textit{La Seine} (floating form), \textit{le Château de Versailles} (inconsistent)}
    {\textit{samedi}, \textit{la Seine}, \textit{le château de Versailles}}
    {In standard French, weekdays and most common nouns stay lowercase.}
    {\textit{Samedi dernier, j'ai visité le château de Versailles.}}
    {5 occurrences}
    {A2 Sample 1 Ex.1; A2 Sample 3 Ex.1; B1 Exemple 1}
\end{enumerate}

\subsection*{Micro-drills}
\begin{enumerate}
    \item Correct: \textit{J'espiere que tu va bien et j'ai visité le Samedi.}
    \item Fill in: \textit{Au cours des \underline{\hspace{1cm}} dernières années, les changements \underline{\hspace{2cm}} sont visibles.}
\end{enumerate}
\textbf{Quick answers:} 1) \textit{J'espère que tu vas bien et j'ai visité samedi.} 2) \textit{vingt, environnementaux}.

\section{Agreement (Gender, Number, Adjectives, Participles)}
\subsection*{Entries}
\begin{enumerate}
    \mistakeentry
    {\textit{mon nouvelle apartement}, \textit{un changement négative}, \textit{les familles du ma mère}}
    {\textit{mon nouvel appartement}, \textit{un changement négatif}, \textit{la famille de ma mère}}
    {Match article/adjective forms with noun gender and number.}
    {\textit{C'est un changement négatif pour la famille de ma mère.}}
    {10 occurrences}
    {A1 Sample 2/4; A2 Sample 3 Ex.1; B1 Exemple 1/2; Writing Practice Day 1-3}

    \mistakeentry
    {\textit{les gens ... est très gentil}, \textit{tout le monde ont}}
    {\textit{les gens ... étaient très gentils}, \textit{tout le monde avait}}
    {Subject-verb agreement is mandatory: plural with \textit{les gens}, singular with \textit{tout le monde}.}
    {\textit{Les gens à qui j'ai parlé étaient très gentils.}}
    {8 occurrences}
    {A2 Sample 2 Ex.1; A2 Sample 3 Ex.1; B1 Exemple 2}

    \mistakeentry
    {Participle/adjective mismatch: \textit{être exposer}, \textit{interésant/intéréssant} in feminine/plural contexts}
    {\textit{être exposé}, \textit{intéressante/intéressants} based on target noun}
    {With \textit{être}, participles/adjectives must agree with the subject.}
    {\textit{Cette expérience a été intéressante et enrichissante.}}
    {6 occurrences}
    {A2 Sample 2 Ex.1; B1 Exemple 1; Writing Practice Day 1-3}
\end{enumerate}

\subsection*{Micro-drills}
\begin{enumerate}
    \item Correct: \textit{Tout le monde ont arrière-plan différent.}
    \item Fill in: \textit{Les personnes à qui j'ai parlé étaient très \underline{\hspace{2cm}}.}
\end{enumerate}
\textbf{Quick answers:} 1) \textit{Tout le monde avait un arrière-plan différent.} 2) \textit{gentilles/gentils}.

\section{Verb Forms and Tenses}
\subsection*{Entries}
\begin{enumerate}
    \mistakeentry
    {\textit{J'ai arrivé}, \textit{Je arriverais} (for a certain planned future)}
    {\textit{Je suis arrivé}, \textit{Je viendrai / J'arriverai}}
    {Many movement verbs take \textit{être} in \textit{passé composé}; use futur simple for scheduled future events.}
    {\textit{Je suis arrivé hier soir et j'arriverai samedi à 16 h.}}
    {7 occurrences}
    {A1 Sample 2/3; A2 Sample 1 Ex.1; A2 Sample 4 Ex.2; A2 Sample 3 Ex.2}

    \mistakeentry
    {\textit{j'espèrais que tu va bien}, \textit{je n'serai pas}}
    {\textit{j'espère que tu vas bien}, \textit{je ne serai pas}}
    {Stabilize core paradigms (present/future) and keep full negation form \textit{ne ... pas}.}
    {\textit{J'espère que tu vas bien et je ne serai pas en retard.}}
    {9 occurrences}
    {A2 Sample 1 Ex.2; A1 Sample 3; Writing Practice Day 1-3; B1 Exemple 2}

    \mistakeentry
    {\textit{après obtenir son diplôme}}
    {\textit{après avoir obtenu mon diplôme}}
    {After \textit{après}, use \textit{avoir/être + participe passé} to express completed prior action.}
    {\textit{Je suis revenu en Chine après avoir obtenu mon diplôme.}}
    {3 occurrences}
    {B1 Exemple 1; A2 Sample 1 Ex.1}
\end{enumerate}

\subsection*{Micro-drills}
\begin{enumerate}
    \item Correct: \textit{J'ai arrivé à Paris et je arriverais demain soir.}
    \item Fill in: \textit{Après avoir \underline{\hspace{2cm}} mes bagages, je \underline{\hspace{2cm}} de l'aéroport.}
\end{enumerate}
\textbf{Quick answers:} 1) \textit{Je suis arrivé à Paris et j'arriverai demain soir.} 2) \textit{récupéré, sortirai} (or \textit{suis sorti} if changing tense context).

\section{Prepositions and Articles}
\subsection*{Entries}
\begin{enumerate}
    \mistakeentry
    {\textit{dans le 2 octobre}, \textit{dans le 20 juillet 2026}}
    {\textit{le 2 octobre}, \textit{le 20 juillet 2026}}
    {For exact dates, use \textit{le + date} directly; do not add \textit{dans}.}
    {\textit{Mon vol arrivera à Paris le 20 juillet 2026.}}
    {6 occurrences}
    {A2 Sample 2 Ex.1; A2 Sample 3 Ex.2; recurring A2 notes}

    \mistakeentry
    {\textit{au Étas-Unis}, \textit{dans Paris}, \textit{en Paris}}
    {\textit{aux États-Unis}, \textit{à Paris}}
    {Use \textit{à} with cities; \textit{aux} with plural countries; respect mandatory contractions.}
    {\textit{Je vais à Paris, puis je pars aux États-Unis.}}
    {8 occurrences}
    {A2 Sample 1 Ex.1; A2 Sample 1 Ex.2; A2 Sample 3 Ex.2; B1 Exemple 1}

    \mistakeentry
    {\textit{de émissions}, \textit{du ma mère}, \textit{à le marché}}
    {\textit{d'émissions}, \textit{de ma mère}, \textit{au marché}}
    {Apply elision/contractions: \textit{de + vowel = d'}, \textit{à + le = au}, \textit{de + le = du}.}
    {\textit{Il y a moins d'émissions et je vais au marché du quartier.}}
    {12 occurrences}
    {A1 Sample 2; A2 Sample 3 Ex.1; B1 Exemple 1/2; Writing Practice Day 1-3}
\end{enumerate}

\subsection*{Micro-drills}
\begin{enumerate}
    \item Correct: \textit{Mon vol est dans le 15 août et je vais à le centre.}
    \item Fill in: \textit{Je réduis le bruit et \underline{\hspace{1cm}} émissions polluantes en ville.}
\end{enumerate}
\textbf{Quick answers:} 1) \textit{Mon vol est le 15 août et je vais au centre.} 2) \textit{d'}.

\section{Relative and Object Pronouns}
\subsection*{Entries}
\begin{enumerate}
    \mistakeentry
    {\textit{que est un palais}, \textit{les gens que j'ai parlé}}
    {\textit{qui est un palais}, \textit{les gens à qui j'ai parlé}}
    {\textit{Qui} = subject, \textit{que} = direct object; with \textit{parler à}, use \textit{à qui}.}
    {\textit{Le château qui est célèbre attire les gens à qui j'ai parlé.}}
    {8 occurrences}
    {A2 Sample 1 Ex.1; A2 Sample 2 Ex.1; B1 Exemple 1}

    \mistakeentry
    {\textit{j'invite tu visites mon pays}}
    {\textit{je t'invite à visiter mon pays}}
    {Object pronouns (\textit{me/te/le...}) come before the conjugated verb.}
    {\textit{Je t'invite à visiter mon pays cet été.}}
    {4 occurrences}
    {A1 Sample 4; A2 Sample 4 Ex.2}

    \mistakeentry
    {\textit{qui forme un réseau de surveillance} (referring to a full clause)}
    {\textit{ce qui crée un réseau de surveillance}}
    {When referring to a whole previous idea, prefer \textit{ce qui/ce que}.}
    {\textit{Les caméras sont partout, ce qui crée un réseau de surveillance.}}
    {3 occurrences}
    {B1 Exemple 1; B1 Exemple 2}
\end{enumerate}

\subsection*{Micro-drills}
\begin{enumerate}
    \item Correct: \textit{Le livre que je parle est intéressant.}
    \item Fill in: \textit{Les personnes \underline{\hspace{1cm}} j'ai parlé étaient chaleureuses.}
\end{enumerate}
\textbf{Quick answers:} 1) \textit{Le livre dont je parle est intéressant.} 2) \textit{à qui}.

\section{Syntax and Word Order}
\subsection*{Entries}
\begin{enumerate}
    \mistakeentry
    {\textit{a changé beaucoup}, \textit{la Chine a été concentrer les ressources}}
    {\textit{a beaucoup changé}, \textit{la Chine s'est concentrée sur les ressources}}
    {Use standard adverb placement and idiomatic verb structures.}
    {\textit{Au cours des dernières années, la Chine s'est concentrée sur ce secteur et a beaucoup changé.}}
    {9 occurrences}
    {B1 Exemple 1; B1 Exemple 2; Writing Practice Day 1-3}

    \mistakeentry
    {Long run-on sequences with weak logical segmentation}
    {Split into shorter clauses/sentences + precise connectors}
    {Avoid stacking clauses without hierarchy; aim for one main idea per sentence.}
    {\textit{L'événement était surprenant. En revanche, je n'ai pas aimé la nourriture.}}
    {10 occurrences}
    {A2 Sample 2 Ex.1; B1 Exemple 1/2; Writing Practice Day 1-3}

    \mistakeentry
    {\textit{Je pense que serais amusante et spectaculaire}}
    {\textit{Je pense que ce sera amusant et spectaculaire}}
    {After \textit{je pense que}, keep a complete clause with an explicit subject.}
    {\textit{Je pense que ce sera une expérience très enrichissante.}}
    {5 occurrences}
    {A1 Sample 4; A2 Sample 3 Ex.1; B1 Exemple 2}
\end{enumerate}

\subsection*{Micro-drills}
\begin{enumerate}
    \item Correct: \textit{Il y a beaucoup des gens la ville est sale il faut nettoyer.}
    \item Fill in: \textit{Je pense que \underline{\hspace{2cm}} utile pour les élèves.}
\end{enumerate}
\textbf{Quick answers:} 1) \textit{Il y a beaucoup de gens. La ville est sale, donc il faut nettoyer.} 2) \textit{ce sera / c'est}.

\section{Lexicon, Collocations, and Register}
\subsection*{Entries}
\begin{enumerate}
    \mistakeentry
    {\textit{jouer des spectacles}, \textit{manger diner}, \textit{terrain du joue}}
    {\textit{présenter des spectacles}, \textit{dîner}, \textit{cour de récréation/terrain de jeu}}
    {Memorize frequent collocations instead of literal translation from English.}
    {\textit{J'ai vu des artistes présenter des spectacles de rue.}}
    {9 occurrences}
    {A2 Sample 2 Ex.1; A2 Sample 3 Ex.1; Writing Practice Day 1-3}

    \mistakeentry
    {\textit{gens} in counted/formal contexts, \textit{recontrer} for a person already known}
    {\textit{personnes} (often better), \textit{retrouver} (if not a first meeting)}
    {Choose words based on nuance, register, and discourse context.}
    {\textit{Il y avait douze personnes, et j'ai retrouvé une amie à la gare.}}
    {6 occurrences}
    {A2 Sample 3 Ex.1/2; A2 notes; B1 Exemple 2}

    \mistakeentry
    {False friends/calques: \textit{proper}, \textit{case}, \textit{forum} (imprecise usage)}
    {\textit{propre}, \textit{cas}, \textit{thème/sujet} depending on context}
    {Check near-English words carefully; meaning/register often differs.}
    {\textit{Le même cas s'applique aux citoyens japonais.}}
    {5 occurrences}
    {B1 Exemple 2; A2 Sample 4 Ex.1; Writing Practice Day 1-3}
\end{enumerate}

\subsection*{Micro-drills}
\begin{enumerate}
    \item Correct: \textit{J'ai joué un spectacle et la ville est proper.}
    \item Fill in: \textit{J'ai vu des artistes \underline{\hspace{2cm}} des performances culturelles.}
\end{enumerate}
\textbf{Quick answers:} 1) \textit{J'ai présenté/assisté à un spectacle et la ville est propre.} 2) \textit{présenter}.

\section{Logical Connectors}
\subsection*{Entries}
\begin{enumerate}
    \mistakeentry
    {\textit{depuis il est trop facile de suivre quelqu'un}}
    {\textit{car / puisque il est trop facile de suivre quelqu'un}}
    {\textit{Depuis} is temporal; for cause, use \textit{car/parce que/puisque}.}
    {\textit{Je suis inquiet, car il est trop facile de suivre quelqu'un.}}
    {4 occurrences}
    {B1 Exemple 1; Writing Practice Day 1-3}

    \mistakeentry
    {\textit{Bien que ..., mais ...}}
    {\textit{Bien que ..., ...} \textbf{or} \textit{Même si ..., ... mais ...} depending on structure}
    {Avoid double concession markers in one standard clause structure.}
    {\textit{Bien que la technologie facilite la vie, elle pose des risques.}}
    {3 occurrences}
    {B1 Exemple 1; B1 Exemple 2}

    \mistakeentry
    {Missing or weak connector calibration across successive ideas}
    {Use \textit{d'abord, ensuite, en revanche, donc, en conclusion} by logical relation}
    {Each connector should encode a precise relation: addition, contrast, consequence, conclusion.}
    {\textit{D'abord..., ensuite..., en revanche..., en conclusion...}}
    {7 occurrences}
    {A1 Sample 4; A2 Sample 2 Ex.1; B1 Exemple 2; Writing Practice Day 1-3}
\end{enumerate}

\subsection*{Micro-drills}
\begin{enumerate}
    \item Correct: \textit{Bien que c'est utile, mais c'est dangereux.}
    \item Fill in with a causal connector: \textit{Je refuse ce choix, \underline{\hspace{1cm}} il est trop risqué.}
\end{enumerate}
\textbf{Quick answers:} 1) \textit{Bien que ce soit utile, c'est dangereux.} 2) \textit{car / parce que / puisque}.

\section{Pre-Writing Checklist}
\begin{enumerate}
    \item Did I check date format (\textit{le 20 juillet}, not \textit{dans le})?
    \item Did I check \textit{à/en/au/aux} and contractions \textit{au/du/d'}?
    \item Did I verify agreement (gender, number, subject-verb)?
    \item Are my relative pronouns correct (\textit{qui/que/dont/à qui})?
    \item Did I avoid English calques and false friends?
    \item Are my connectors precise and non-redundant?
\end{enumerate}

\end{document}
